\section*{Chapitre \ref{ch:b1:fluid-statics} --- Statique des fluides}

\begin{solution}{exo:b1:fluid-statics:1}
$P = P_0 + \rho gh$~: \qty{2.0e5}{Pa} (\qty{2.0}{bar}), \qty{1.1e6}{Pa}
(\qty{11}{bar}), \qty{1.0e7}{Pa} (\qty{101}{bar}). À \qty{300}{m}~: $P =
\qty{3.1e6}{Pa}$, force \qty{3.1}{MN} sur \qty{1}{m^2} (résultante, contre
\qty{1}{atm} à l'intérieur~: \qty{3.0}{MN}).
\end{solution}

\begin{solution}{exo:b1:fluid-statics:2}
$h = P/\rho g$~: mercure \qty{0.760}{m}~; eau \qty{10.3}{m}. Une pompe ne peut
au mieux que faire le vide au-dessus de l'eau~: l'atmosphère pousse alors la
colonne de \qty{10.3}{m}, pas davantage.
\end{solution}

\begin{solution}{exo:b1:fluid-statics:3}
$F_2 = F_1S_2/S_1 = \qty{10}{kN}$ (une tonne)~; $d_2 = d_1S_1/S_2 = \qty{2.0}{mm}$
(même volume déplacé, même travail).
\end{solution}

\begin{solution}{exo:b1:fluid-statics:4}
Même pression au niveau de l'interface huile--eau dans les deux branches~: $\rho_og
h_o = \rho_wgh_w$, soit $h_w = 0.85 \times 12 = \qty{10.2}{cm}$ d'eau au-dessus
de ce niveau dans l'autre branche~; la surface de l'huile est $12 - 10.2 = \qty{1.8}{cm}$
plus haut.
\end{solution}

\begin{solution}{exo:b1:fluid-statics:5}
$H = RT/Mg = 8.314 \times 273/(0.029 \times 9.81) = \qty{8.0}{km}$~; $P/P_0 =
\eu^{-z/H}$~: $0.54$ à \qty{5}{km}, $0.33$ à \qty{8848}{m}~; la moitié à $H\ln2 =
\qty{5.5}{km}$ --- les valeurs utilisées au \cref{ch:b1:kinetic-theory}.
\end{solution}

\begin{solution}{exo:b1:fluid-statics:6}
$F = \tfrac12\rho gh^2L = \qty{1.96e8}{N}$~; \omterm{prop:b1:fluid-statics:wall}{centre de poussée} à \qty{13.3}{m}
de profondeur, soit \qty{6.7}{m} au-dessus du pied~; moment par rapport au pied $F \times
6.7 = \qty{1.3e9}{N.m}$. Courber le barrage vers l'eau transforme la poussée
en compression de l'arc, transmise aux flancs de la vallée.
\end{solution}

\begin{solution}{exo:b1:fluid-statics:7}
$917/1025 = 89\%$~; bûche $60\%$. Radeau~: poussée à immersion totale $1000 \times
2.0 \times 9.81 = \qty{19.6}{kN}$, poids propre \qty{11.8}{kN}~: \qty{7.8}{kN}
disponibles, soit dix personnes\dots{} la dixième les pieds dans l'eau~: neuf pour rester au sec
(\qty{800}{kg}).
\end{solution}

\begin{solution}{exo:b1:fluid-statics:8}
Volume immergé $m/\rho$~: eau \qty{20}{cm^3}~; saumure $20/1.1 = \qty{18.2}{cm^3}$,
soit $\qty{1.8}{cm^3}$ de moins, c'est-à-dire $1.8/0.50 = \qty{3.6}{cm}$ de tige en plus~: \qty{7.6}{cm}
émergent. Sensibilité $\approx m/(\rho^2S) \approx \qty{40}{cm}$ par unité de densité
(\qty{3.6}{cm} pour $0.1$).
\end{solution}

\begin{solution}{exo:b1:fluid-statics:9}
Volume déplacé $2800/1.025 = \qty{2732}{m^3}$~: \qty{268}{m^3} ($9\%$)
émergent. Flottabilité nulle~: masse $1025 \times 3000 = \qty{3075}{t}$, donc \qty{275}{t}
à embarquer. En eau douce, la poussée tombe à \qty{3000}{t}~: \qty{75}{t} de trop --- il coule
si l'on n'en chasse pas \qty{75}{t}.
\end{solution}

\begin{solution}{exo:b1:fluid-statics:10}
Air déplacé~: $10^{-2} \times 1.20 = \qty{12.0}{g}$~; hélium intérieur $12.0 \times
4/29 = \qty{1.7}{g}$~; caoutchouc \qty{3.0}{g}~: force ascensionnelle nette \qty{7.3}{g} --- sept mètres
de ficelle légère.
\end{solution}

\begin{solution}{exo:b1:fluid-statics:11}
Moment $\int_0^h\rho gx\cdot xL\,\dd x = \rho gLh^3/3$, force $\rho gLh^2/2$~:
\omterm{def:b1:angular-momentum:moment}{bras de levier} $2h/3$. De $h_1$ à $h_2$~: moment $\rho gL(h_2^3 - h_1^3)/3$, force
$\rho gL(h_2^2 - h_1^2)/2$~: profondeur $\frac23(h_2^3 - h_1^3)/(h_2^2 - h_1^2)$.
\end{solution}

\begin{solution}{exo:b1:fluid-statics:12}
Une immersion supplémentaire $x$ ajoute la poussée $\rho gSx$ vers le haut~: $m\ddot x = -\rho gSx$,
$\omega_0^2 = \rho gS/m$, $T = 2\pi\sqrt{m/\rho gS} = 2\pi\sqrt{0.020/(1000 \times
9.81 \times 5\times10^{-5})} = \qty{1.3}{s}$.
\end{solution}

\begin{solution}{pb:b1:fluid-statics:1}
\textbf{1.} $P = 1.013\times10^5 + 1025 \times 9.81 \times 300 = \qty{3.12e6}{Pa}
= \qty{31}{bar}$.

\textbf{2.} $S = \pi(0.40)^2 = \qty{0.503}{m^2}$~: force résultante $(P - P_0)S =
\qty{1.52e6}{N}$ --- 155 tonnes, un camion chargé sur un panneau.

\textbf{3.} La quille est \qty{10}{m} plus bas~: $+\rho g \times 10 = \qty{1.0e5}{Pa}$ de plus,
soit $3\%$~: la coque est chargée presque uniformément.

\textbf{4.} $\Delta\rho/\rho = \chi_T\Delta P = 4.6\times10^{-10} \times 3.0\times10^6
= 0.14\%$~: négligeable pour la pression (quelques kilopascals sur trois
millions).

\textbf{5.} Trente bars au-dehors, un au-dedans~: la coque est une enceinte
sous pression chargée vers l'intérieur~; un cylindre (ou une sphère) transforme
la pression en compression le long de sa paroi, ce que l'acier supporte bien,
alors qu'une plaque plane fléchirait.

\textbf{6.} $\Delta P\,R/e = 3.0\times10^6 \times 5.0/0.030 = \qty{5.0e8}{Pa}
= \qty{500}{MPa}$~: proche de la limite d'élasticité --- \qty{300}{m} est près
de la limite pour cette coque, et les bâtiments plus profonds demandent un
acier plus épais ou plus résistant (ou du titane).

\textbf{7.} Volume déplacé $m_0/\rho = 2800/1.025 = \qty{2732}{m^3}$~: \qty{268}{m^3}
émergent ($9\%$).

\textbf{8.} Flottabilité nulle~: $m = \rho V = \qty{3075}{t}$, donc \qty{275}{t} à embarquer.

\textbf{9.} La poussée diminue de $0.1\%$ de $\rho Vg$~: $0.001 \times 3075 \times 9.81
\times10^3 = \qty{30}{kN}$ vers le bas (trois tonnes). Plus profond $\Rightarrow$ plus
comprimé $\Rightarrow$ plus lourd~: équilibre instable --- un sous-marin doit corriger
son assiette en permanence (et une coque d'acier se comprime moins que l'eau,
ce qui aide~; une coque trop souple coulerait sans retour).

\textbf{10.} Poussée $1000 \times 3000 \times 9.81 = \qty{29.4}{MN}$ contre un poids
$3075 \times 9.81 \times 10^3 = \qty{30.2}{MN}$~: \qty{0.74}{MN} vers le bas
(\qty{75}{t})~; il faut en chasser \qty{75}{t}.

\textbf{11.} Aussitôt \qty{20}{t} trop léger~: embarquer \qty{20}{t} d'eau
(\qty{19.5}{m^3}) dans les caisses de compensation.

\textbf{12.} L'eau ne peut être chassée que par quelque chose à une pression
supérieure à celle de la mer au-dehors~; l'air comprimé des bouteilles l'expulse
--- mais seulement tant que la pression des bouteilles dépasse la pression
ambiante à cette profondeur.

\textbf{13.} $PV$ constant~: \qty{275}{m^3} à \qty{31}{bar} demandent $275 \times 31
= \qty{8500}{m^3}$ d'air à \qty{1}{atm} --- stockés à \qty{200}{bar} dans
\qty{43}{m^3} de bouteilles.

\textbf{14.} \qty{2.0}{bar} à \qty{10}{m}~; \qty{4.0}{bar} à \qty{30}{m}.

\textbf{15.} $PV$ constant~: $6.0/4.0 = \qty{1.5}{L}$ --- la cage thoracique est
écrasée~; les apnéistes le sentent.

\textbf{16.} L'air pris à \qty{4}{bar} se dilate quatre fois~: \qty{24}{L} dans des
poumons qui en contiennent \qty{6}{L}~: rupture. Ne jamais bloquer sa respiration
à la remontée~; expirer continûment.

\textbf{17.} $\Delta P = \rho g \times 1.0 = \qty{1.0e4}{Pa}$~; force $\qty{1.0e4}{Pa}
\times 0.05 = \qty{500}{N}$ sur la poitrine~: les muscles respiratoires ne
soulèvent pas cinquante kilogrammes --- un tuba ne sert que dans les premiers décimètres.

\textbf{18.} $12 \times 200 = \qty{2400}{L}$ à \qty{1}{bar}, soit \qty{600}{L}
à \qty{4}{bar}~: \qty{30}{min}.

\textbf{19.} $V = 10^8/1025 = \qty{97600}{m^3}$.

\textbf{20.} En eau douce, il doit déplacer $10^5\ \unit{m^3}$~: \qty{2400}{m^3}
de plus, soit $2400/8000 = \qty{0.30}{m}$ d'enfoncement en plus.

\textbf{21.} $10^4/1.025 = \qty{9756}{m^3}$ de plus~: \qty{1.2}{m} plus bas.

\textbf{22.} À vide, le $G$ du navire est haut et peu de coque est immergée~:
le métacentre peut passer sous $G$ et le navire chavirer dans la houle. De
l'eau de mer basse dans la coque abaisse $G$ et augmente le tirant d'eau.

\textbf{23.} $\rho gh = \qty{2.0e5}{Pa}$, soit \qty{0.2}{MN/m^2}~: quinze fois moins
que les \qty{3}{MN/m^2} du sous-marin.

\textbf{24.} Non~: la pression à la profondeur du sous-marin vaut $P_0 + \rho gh$,
fixée par la seule profondeur. Le poids du pétrolier est porté par l'eau qu'il
déplace, ce qui élève partout le niveau de la mer d'une quantité immesurable~;
aucune charge supplémentaire n'atteint le sous-marin.

\textbf{25.} $\dd P/\dd z = -\rho g$ --- d'où $P = P_0 + \rho gh$ pour chaque
profondeur, chaque force et chaque volume pulmonaire de ce problème --- et le
théorème d'Archimède, qui n'est que cette relation intégrée sur une surface fermée.
\end{solution}
